\section{A spherical sign estimate}\label{sec:spherical}

The estimate in this section separates the constant part of a spherical
function from its harmonics of degree at least four.  The degree-two part
will later be removed by a linear change of coordinates.  Throughout,
$p=n-1$, surface measure $\sigma$ has total mass one, and inner products
and norms of functions are those of $L^2(\sigma)$.

\subsection{Harmonics and the spectral gap}

Let $\mathcal H_l$ be the restrictions to $\Sph$ of homogeneous harmonic
polynomials of degree $l$, and put
\[
 \lambda_l=l(l+n-2)=l(l+p-1).
\]
For an even function, let $Q$ denote orthogonal projection onto the
closed sum of $\mathcal H_4,\mathcal H_6,\ldots$.

\begin{lemma}\label{lem:spherical-harmonics}
The spaces $\mathcal H_l$ form an orthogonal decomposition of
$L^2(\sigma)$, and $\lapS Y=-\lambda_lY$ for $Y\in\mathcal H_l$.
Every even $H\in C^1(\Sph)$ satisfies
\begin{equation}\label{eq:spectral-gap}
 \lambda_2\Var(H)+(\lambda_4-\lambda_2)\norm{QH}_2^2
 \le \E\abs{\gradS H}^2,
 \qquad \lambda_2=2n,\quad \lambda_4=4(n+2).
\end{equation}
\end{lemma}

\begin{proof}
The polar-coordinate formula for the Euclidean Laplacian gives
\[
 \Delta_{\R^n}\big(r^lY(u)\big)
 =r^{l-2}\big(l(l+n-2)Y(u)+\lapS Y(u)\big).
\]
This proves the eigenvalue formula.  Integration by parts on the sphere
then proves orthogonality of spaces with distinct degrees.

We give the completeness argument explicitly.  If $Y_l$ is homogeneous
and harmonic, direct differentiation yields
\[
 \Delta_{\R^n}\big(\abs{x}^{2j}Y_l(x)\big)
 =2j(2l+2j+n-2)\abs{x}^{2j-2}Y_l(x).
\]
Induction on the degree now decomposes every homogeneous polynomial
into terms $\abs{x}^{2j}Y_l(x)$.  Indeed, decompose its Laplacian by the
induction hypothesis and use the displayed identity to find a sum of
such terms with the same Laplacian.  The difference is harmonic.
Restricting to the unit sphere shows that every polynomial restriction
belongs to the span of the $\mathcal H_l$.

Polynomial restrictions uniformly approximate continuous functions on
$\Sph$.  To see this without an approximation theorem, set
\[
 K_j(u,v)=
 \frac{((1+u\cdot v)/2)^j}
 {\displaystyle\int_{\Sph}((1+u\cdot w)/2)^j\,d\sigma(w)}.
\]
The denominator is positive and independent of $u$, so
$T_jH(u)=\int K_j(u,v)H(v)\,d\sigma(v)$ is a polynomial in $u$.
For $0<\delta<2$, let $m_\delta>0$ be the measure of the cap
$\{v:\abs{u-v}<\delta/2\}$.  Since
$(1+u\cdot v)/2=1-\abs{u-v}^2/4$,
\[
 \int_{\abs{u-v}\ge\delta}K_j(u,v)\,d\sigma(v)
 \le \frac1{m_\delta}
 \left(\frac{1-\delta^2/4}{1-\delta^2/16}\right)^j.
\]
Uniform continuity of $H$ therefore gives $T_jH\to H$ uniformly.
Continuous functions are dense in $L^2(\sigma)$: approximate the sets
in a simple function from inside by compact sets and from outside by
open sets, then interpolate their indicators continuously using
distance to these sets.  The squared error is bounded by the measure
of the intervening sets.  This proves completeness.

Choose real orthonormal bases $Y_{l,a}$ of the spaces $\mathcal H_l$,
with $Y_{0,1}=1$.  Integration by parts shows that the tangent fields
$\gradS Y_{l,a}/\sqrt{\lambda_l}$, for $l\ge1$, are orthonormal.
Bessel's inequality, which follows by expanding the squared distance
from any finite orthogonal sum, consequently gives
\[
 \E\abs{\gradS H}^2
 \ge \sum_{l\ge1,a}
 \left|\E\,\gradS H\cdot
       \frac{\gradS Y_{l,a}}{\sqrt{\lambda_l}}\right|^2
 =\sum_{l\ge1,a}\lambda_l\left|\E(HY_{l,a})\right|^2.
\]
No differentiation of an infinite series is needed.  If $H$ is even,
its odd-degree coefficients vanish.  Completeness identifies its
variance with the sum of the squared nonconstant coefficients and
$\norm{QH}_2^2$ with the corresponding sum over even $l\ge4$.
Using $\lambda_l\ge\lambda_2$ for even $l\ge2$ and
$\lambda_l\ge\lambda_4$ for $l\ge4$ proves
\eqref{eq:spectral-gap}.
\end{proof}

\subsection{The sign-test operator}

For $k>0$ and an even differentiable function $f$ on $\Sph$, define
\begin{equation}\label{eq:vector-field}
 X_f(u)=f(u)u+k^{-1}\gradS f(u).
\end{equation}
For even $C^1$ functions $f_1,f_2$, set
\begin{equation}\label{eq:sign-form}
 \mathcal B(f_1,f_2)
 =\frac1{b_n}\iint_{\Sph\times\Sph}
   \sgn(u\cdot v)\,X_{f_1}(u)\cdot X_{f_2}(v)
   \,d\sigma(u)d\sigma(v).
\end{equation}
The value assigned to $\sgn(0)$ is immaterial. Since
$|a|\ge\sgn(u\cdot v)a$ for every real $a$, the form gives the lower bound
\[
 \iint |X_{f_1}(u)\cdot X_{f_2}(v)|\,d\sigma(u)d\sigma(v)
 \ge b_n\mathcal B(f_1,f_2).
\]
We estimate this lower bound using two spherical transforms.
Write $\sigma_{u^\perp}$
for uniform probability on the equator $u^\perp\cap\Sph$, and define
\[
 (Cf)(u)=\frac1{b_n}\int_{\Sph}\abs{u\cdot v}f(v)\,d\sigma(v),
 \qquad
 (Pf)(u)=\int_{u^\perp\cap\Sph}f(v)\,d\sigma_{u^\perp}(v).
\]
These are the cosine and Funk transforms, here normalized to preserve
constants; compare \cite[Section~2, especially (2.3), (2.5)--(2.6)]
{OlafssonPasqualeRubin2013}. The local multiplier calculation below
uses this probability-measure convention throughout.
Both are self-adjoint contractions on $L^2(\sigma)$ and preserve
constants.  For $C$ this follows from symmetry of its nonnegative
kernel and Jensen's inequality.  For $P$, average first over the band
$\abs{u\cdot v}<\epsilon$, dividing by its measure.  The resulting
kernel is symmetric and preserves constants.  Spherical coordinates
show that, for continuous $f$, these band averages tend to $Pf$;
Jensen's inequality and passage to the limit give the same conclusions
for $P$, followed by extension to $L^2(\sigma)$.

\begin{lemma}\label{lem:sign-operator}
For even $C^1$ functions,
\begin{equation}\label{eq:sign-operator}
 \mathcal B(f_1,f_2)
 =\left\langle f_1,
   \left[C+\left(\frac{2p}{k}+\frac{p^2}{k^2}\right)(C-P)\right]
   f_2\right\rangle.
\end{equation}
\end{lemma}

\begin{proof}
Let $d_0$ be the density at zero of $u\cdot v$ for uniform $u$ and fixed
$v\in\Sph$.  Equivalently, $d_0$ is the ratio of the surface areas of
$S^{n-2}$ and $S^{n-1}$.  For a fixed ambient vector $w$, its tangential
projection has divergence $-p\,u\cdot w$.  Integrating separately on
the two hemispheres gives
\begin{align}
 &\int_{\Sph}\sgn(u\cdot v)\,\gradS f(u)\cdot w\,d\sigma(u)
 \notag\\
 &\quad=p\int_{\Sph}\sgn(u\cdot v)f(u)\,u\cdot w\,d\sigma(u)
 -2d_0\int_{v^\perp\cap\Sph}f(u)\,v\cdot w
       \,d\sigma_{v^\perp}(u).
 \label{eq:hemisphere-ibp}
\end{align}
The outward conormals on the equator are $-v$ and $v$, respectively;
the sign change between the hemispheres makes their contributions
add with the displayed negative sign.  Taking $f=1$ and $w=v$ yields
\[
 2d_0=p b_n.
\]

The radial--radial term in \eqref{eq:sign-form}, before division by
$b_n$, is $b_n\langle f_1,Cf_2\rangle$.  Apply
\eqref{eq:hemisphere-ibp} with $w=v$ and then integrate against
$f_2(v)$.  Each of the two mixed radial--gradient terms, before its
factor $k^{-1}$, equals
\[
 p b_n\langle f_1,(C-P)f_2\rangle.
\]
For the gradient--gradient term, first integrate in $u$ with
$w=\gradS f_2(v)$.  Its equator term vanishes because
$v\cdot\gradS f_2(v)=0$.  What remains is
\[
 p\iint f_1(u)\sgn(u\cdot v)\,
      u\cdot\gradS f_2(v)\,d\sigma(u)d\sigma(v).
\]
A second application of \eqref{eq:hemisphere-ibp}, now in $v$ with
$w=u$, makes this $p^2b_n\langle f_1,(C-P)f_2\rangle$.
Adding the four terms proves \eqref{eq:sign-operator}.
\end{proof}

\subsection{Multipliers and passage to norms}

\begin{lemma}\label{lem:spherical-multipliers}
On $\mathcal H_{2j}$ the operators $P$ and $C$ are scalar multipliers,
with
\begin{equation}\label{eq:transform-multipliers}
 P_{2j}=(-1)^j
 \frac{1\cdot3\cdots(2j-1)}{p(p+2)\cdots(p+2j-2)},
 \qquad
 C_{2j}=\frac{pP_{2j}}{p-\lambda_{2j}},
\end{equation}
where empty products equal one.  In particular,
\begin{equation}\label{eq:cosine-recurrence}
 \begin{gathered}
 C_0=1,\qquad C_2=\frac1{p+2},\qquad
 C_4=-\frac1{(p+2)(p+4)},\\[3pt]
 \frac{C_{l+2}}{C_l}=-\frac{l-1}{l+p+2}\quad(l\ge2\text{ even}).
 \end{gathered}
\end{equation}
\end{lemma}

\begin{proof}
Fix $u\in\Sph$ and a homogeneous harmonic polynomial $Y_{2j}$.
Let $Z$ be a standard Gaussian in the $p$-dimensional space $u^\perp$.
Its direction is uniform on the equator and independent of its radius,
so polar integration gives
\[
 \mathbb E Y_{2j}(Z)
 =p(p+2)\cdots(p+2j-2)\,(PY_{2j})(u).
\]
The radial moment here follows by successive integration by parts in
$\int_0^\infty r^{p+2j-1}e^{-r^2/2}\,dr$.
The one-dimensional Gaussian moment recurrence
$\mathbb E Z_1^{2a}=(2a-1)\mathbb E Z_1^{2a-2}$, together with
independence of coordinates, also gives
\[
 \mathbb E Y_{2j}(Z)
 =\frac{(\Delta_{u^\perp})^jY_{2j}(0)}{2^j j!}.
\]
This identity follows term by term on monomials; monomials with an odd
exponent have both sides zero, and the other coefficients are exactly
the multinomial coefficients in the power of the Laplacian.
Harmonicity and commutation of constant-coefficient derivatives imply
\[
 (\Delta_{u^\perp})^jY_{2j}
   =(-1)^j\partial_u^{2j}Y_{2j},
 \qquad
 \partial_u^{2j}Y_{2j}(0)=(2j)!Y_{2j}(u).
\]
Since $(2j)!/(2^j j!)=1\cdot3\cdots(2j-1)$, the formula for $P_{2j}$
follows.

To obtain $C$, integrate by parts on the two hemispheres in the
variable $v$.  Off the equator,
$\lapS\abs{u\cdot v}=-p\abs{u\cdot v}$; the jump of its normal
derivative contributes $2d_0$ times equator averaging.  Thus, for any
smooth $Y$,
\[
 \int_{\Sph}\abs{u\cdot v}\lapS Y(v)\,d\sigma(v)
 =-p\int_{\Sph}\abs{u\cdot v}Y(v)\,d\sigma(v)
   +2d_0(PY)(u).
\]
For $Y\in\mathcal H_l$, division by $b_n$ gives
$-\lambda_l CY=p(PY-CY)$.  The denominator
$p-\lambda_l=-(l-1)(l+p)$ is nonzero for every even $l$, proving the
second formula in \eqref{eq:transform-multipliers}.  Its first three
values give those in \eqref{eq:cosine-recurrence}; division of successive
values, using $P_{l+2}/P_l=-(l+1)/(l+p)$, gives the recurrence.
\end{proof}

Set
\begin{equation}\label{eq:tail-constant}
 c_{n,k}=
 \frac{1+\left(2/k+(n-1)/k^2\right)4(n+2)}{(n+1)(n+3)}.
\end{equation}

\begin{proposition}\label{prop:sign-estimate}
Let $k>0$ and let $f_1,f_2$ be even $C^1$ functions on $\Sph$.
If $f_1$ has zero projection onto $\mathcal H_2$, then
\begin{equation}\label{eq:sign-bound}
 \mathcal B(f_1,f_2)
 \ge (\E f_1)(\E f_2)
       -c_{n,k}\norm{Qf_1}_2\norm{Qf_2}_2.
\end{equation}
The operator identity \eqref{eq:sign-operator} and this estimate also
hold when $f_i=g_i^k$ for arbitrary norms $g_i$ on $\R^n$, with the
vector fields defined almost everywhere.  In that case
\begin{equation}\label{eq:norm-vector-field}
 X_{g_i^k}(u)=g_i(u)^{k-1}\nabla g_i(u).
\end{equation}
\end{proposition}

\begin{proof}
By Lemma~\ref{lem:spherical-multipliers}, the operator in
\eqref{eq:sign-operator} has multiplier one on constants and, on degree
$l$, multiplier
\[
 \mu_l=C_l\left[1+\left(\frac2k+\frac p{k^2}\right)\lambda_l\right].
\]
Indeed, $C_l-P_l=(\lambda_l/p)C_l$.
Put $a=2/k+p/k^2>0$.  For even $l\ge4$,
\[
 (l+p+2)\lambda_l-(l-1)\lambda_{l+2}
 =(p-1)l^2+(p^2-1)l+2(p+1)>0.
\]
Since also $l-1<l+p+2$, it follows that
\[
 \frac{\abs{\mu_{l+2}}}{\abs{\mu_l}}
 =\frac{l-1}{l+p+2}\,
   \frac{1+a\lambda_{l+2}}{1+a\lambda_l}\le1.
\]
Consequently $\abs{\mu_l}\le\abs{\mu_4}=c_{n,k}$ on all even
$l\ge4$.  The constant term of the bilinear form is
$(\E f_1)(\E f_2)$, and its degree-two term vanishes by hypothesis.
Cauchy--Schwarz on the remaining harmonic coefficients proves
\eqref{eq:sign-bound}.  This expansion is justified by the completeness
in Lemma~\ref{lem:spherical-harmonics} and the boundedness of $C$ and
$P$.

We finish by proving the passage to norms, retaining the details
needed when their unit balls have corners.  Let $g$ be a norm and let
$g_\epsilon=g*\rho_\epsilon$ on $\R^n$, where $\rho_\epsilon$ is a
smooth, even, nonnegative function of integral one, supported in the
ball of radius $\epsilon$.  The functions $g_\epsilon$ are smooth,
even and convex.  If $L$ is the Euclidean Lipschitz constant of $g$,
then
\[
 \abs{g_\epsilon-g}\le L\epsilon,
 \qquad \abs{\nabla g_\epsilon}\le L.
\]
At a differentiability point $x$ of $g$, these gradients converge to
$\nabla g(x)$.  Indeed, every convergent subsequence of the bounded
gradients has a limit $a$, and convexity gives, on passing to the
limit,
\[
 g(y)\ge g(x)+a\cdot(y-x)\qquad(y\in\R^n).
\]
Taking $y=x\pm th$ and then $t\downarrow0$ forces
$a\cdot h=\nabla g(x)\cdot h$ for every $h$, so the only possible
limit is $\nabla g(x)$.

Lemma~\ref{lem:gauge-variation} supplies differentiability of $g$ at
almost every spherical direction.  Equivalently, one may use its
almost-everywhere differentiability in $\R^n$ and homogeneity: the
exceptional set is a union of rays, so polar integration makes its
intersection with the sphere null.  On the sphere $g$ is bounded away
from zero.  Thus the restrictions
$f_\epsilon=g_\epsilon^k|_{\Sph}$ converge uniformly to $g^k$, and
\[
 \gradS f_\epsilon(u)
 =k g_\epsilon(u)^{k-1}
   (\Id-uu^t)\nabla g_\epsilon(u)
 \longrightarrow
 k g(u)^{k-1}(\Id-uu^t)\nabla g(u)
\]
almost everywhere and in $L^2(\sigma)$, by bounded convergence.
Euler's identity $u\cdot\nabla g(u)=g(u)$, obtained by differentiating
$g(tu)=tg(u)$, identifies the limiting vector field with
\eqref{eq:norm-vector-field}.

The approximants $g_\epsilon$ need not be homogeneous; only their
restrictions and tangential gradients are used here.  The fields
$X_{f_\epsilon}$ converge in $L^2$, so the bounded sign kernel in
\eqref{eq:sign-form} permits passage to the limit.  The right side of
\eqref{eq:sign-operator} also converges, because $C$ and $P$ are
bounded on $L^2$.  This proves the operator identity for norm powers.
Finally apply its harmonic multiplier estimate directly to the
limiting functions.  The approximants themselves need not have
vanishing degree-two part: that hypothesis is used only for the
limiting $f_1$.  This proves \eqref{eq:sign-bound} in the stated
nonsmooth setting.
\end{proof}
